\chapter{Diskussion}

\section{Beantwortung der Forschungsfragen}

\subsection{Forschungsfrage 1}

\subsection{Forschungsfrage 2}

\subsection{Forschungsfrage 3}

\section{Übergreifende Diskussion}

    % Was bedeuten die Ergebnisse insgesamt?
    % Unterschiede GPT/Gemini
    % Datenbindung, Zuverlässigkeit, Halluzinationen
    % Einordnung in Literatur / Theorie

\section{Limitationen}

    % Datensatz
    % nur zwei Modelle
    % Anzahl Wiederholungen
    % Bewertungsentscheidungen
    % technische Einschränkungen
    % Generalisierbarkeit

% =====================================================================
% DISKUSSION -- STUB: Werkzeuggestuetzt vs. nativ (noch auszuformulieren)
% ---------------------------------------------------------------------
% - Beobachtung: ChatGPT-Weboberflaeche mit hochgeladener manipulierter CSV
%   rechnet die korrekten 7,9 % aus; der API-Lauf (Luna) liefert die
%   memorisierten 30,5 %.
% - Haupthebel: Weboberflaeche nutzt den Code-Interpreter (Advanced Data
%   Analysis), fuehrt echten Python-Code (groupby('OverTime')['Attrition']
%   .mean()) auf der Datei aus -> deterministisch aus den Daten, kann nicht
%   memorisieren.
% - Gegenstueck API: Datensatz als In-Context-Text, kein Werkzeug. Modell
%   muesste ueber 1470 Zeilen selbst aggregieren, tut das nicht zuverlaessig
%   (E1-Teilmengenbefund) und fuellt die Luecke mit der memorisierten Statistik.
% - Sekundaerfaktoren, gleiche Richtung: anderes Produkt/Modellversion,
%   werkzeugfoerdernder System-Prompt, Datei als Anhang statt Prompt-Text,
%   evtl. Reasoning-Modus. Haupthebel bleibt die Werkzeugnutzung.
% - Korroboration: Web-Ergebnis 7,9 % bestaetigt, dass der manipulierte
%   Datensatz sauber ist und die API-30,5 % Memorisierung sind, kein Datenfehler.
% - Geltungsbereich praezisieren: Befund gilt fuer die native, werkzeuglose
%   In-Context-Verarbeitung, nicht fuer werkzeuggestuetzte EDA-Systeme, die
%   intern Code ausfuehren. Praezise: ohne externes Rechenwerkzeug geben die
%   Modelle die memorisierte Beziehung wieder; mit Code-Werkzeug verschiebt
%   sich das Bild.
% - Ausblick: Grenze zwischen nativer Modellzuverlaessigkeit und werkzeug-
%   gestuetzter Berechnung. Die Experimente isolieren bewusst die native Seite.
% - Merker: Web-Beobachtung nicht als kontrolliertes Ergebnis fuehren, nur als
%   Illustration in der Diskussion. Noch gegenpruefen, ob die Weboberflaeche
%   einen sichtbaren Code-/Analyseschritt gezeigt hat.
% =====================================================================